摘要
近年来，随着智慧社区建设持续推进，无人机巡检凭借机动性强、空间可达性高、作业效率突出的优势，逐步成为社区安全治理的核心手段，
而多无人机与物业人员的空地协同联合巡检，是打通巡检全流程闭环、提升社区治理效能的关键环节。
本文针对智慧社区多无人机 - 物业人员联合巡检优化问题，基于两阶段闭环优化的核心思想，
将整体问题拆解为多无人机协同巡检路径与停留时间优化、空地联合巡检闭环优化两个核心子问题，
以最小化巡检总完成时间为核心目标，兼顾总能耗控制、高优先级点提前完成与多机负载均衡，
构建了完整的多目标优化模型体系，设计了分阶段的启发式求解算法与全局优化策略，为智慧社区巡检的资源配置、
阈值设计与方案规划提供了系统的定量决策支持。
针对多无人机协同巡检路径与停留时间优化问题，本文明确了目标点分配、任务批次规划、访问顺序设计与累计悬停时间分配四大核心决策变量，
在满足基础巡检要求、单次飞行能耗上限与运营时长约束的前提下，构建了以最小化无人机巡检阶段完成时间为核心，
同时兼顾总飞行能耗、高优先级点提前完成度与多机负载均衡的多目标优化模型，
将其转化为带悬停时间分配与电量约束的多旅行商扩展问题开展求解。算法设计方面，
本文提出构造式贪心算法与层次滚动重构算法两套求解方案，结合 K-means 聚类实现初始任务预分组与起点优选，
通过插入、交换、跨任务重分配与 2-opt 等局部搜索算子实现解的迭代优化。
实验结果表明，两套算法在不同无人机数量配置下呈现差异化优势，构造式贪心算法在 1-3 架无人机配置下表现更优，
层次滚动重构算法在 4 架无人机配置下效果更佳；随着无人机数量从 1 增加至 4，巡检阶段完成时间压缩幅度超 75%，
所有方案均能严格满足能耗与运营时长约束，为无人机配置规模选择提供了完整的量化依据。
针对多无人机与物业人员联合巡检闭环优化问题，本文在第一阶段空中巡检优化的基础上，
引入无人机直接确认判定规则，建立了空中巡检与地面复核联动的联合优化模型，
核心解决了无人机悬停时长追加与人工复核工作量缩减之间的边际权衡问题，以最小化总闭环完成时间为核心目标，
同时兼顾高优先级点直接确认比例提升与人工复核节点数量减少。算法设计方面，本文构建了基于边际收益的状态评分函数，
设计了直接确认判定与地面复核路径重算一体化的联合求解器，将物业地面复核路径优化转化为经典旅行商问题进行求解，
同时引入跨任务单点交换算子实现解的局部增强。实验结果表明，相较于仅满足基础巡检要求的基线方案，
本文设计的联合优化方案在 4 架无人机配置下可将总闭环时间压缩超 20%，直接确认节点从 0 个提升至 7 个，
大幅降低物业复核工作量；在物业人员必须待无人机全部返航后统一出发的严格串行规则下，
边际收益优先的串行策略显著优于盲目耗尽空中资源的贪心策略，验证了核心优化逻辑的合理性与有效性。
为进一步挖掘联合优化方案的潜力，验证算法各组件的有效性与结果的稳定性，本文设计自适应大邻域搜索 ALNS 算法开展全局联合搜索，同时完成了 K-means 起点选取消融实验、直接确认阈值倍率扰动实验、不同串行策略对照实验等多组对照验证，并构建了包含总闭环时间、优先级加权直接确认率、人工服务时间占比的多维度补充评价体系。实验结果表明，ALNS 算法在 4 架无人机配置下可实现最为显著的优化效果，能通过重构任务航线进一步压缩总闭环时间，将直接确认节点提升至 12 个；全配置的阈值灵敏度实验明确了 0.8 倍为最优阈值倍率，为社区巡检的阈值设计提供了定量依据；多组消融与对照实验，完整验证了本文所设计算法各组件的有效性与优化结果的稳定性。
本文完整构建了智慧社区多无人机 - 物业人员联合巡检的两阶段闭环优化框架，完整覆盖从模型构建、算法设计、实验验证到决策支持的全流程研究，明确了不同无人机配置下的最优算法选择与方案设计。综合全配置的实验结果，4 架无人机是兼顾巡检效率、资源投入与闭环效果的最优配置规模，在强调稳健合规与可复现性的场景下，推荐采用 4 架无人机配置下的边际收益优先串行联合求解方案，也可根据实际需求叠加局部增强、阈值优化等手段进一步提升效果。本文所建立的模型与算法体系，完整解决了赛题提出的各项优化问题，也为智慧社区巡检的资源配置、阈值设计与路径规划提供了可落地的定量决策支持，具备良好的可复现性与可扩展性。
关键词：多无人机巡检；联合闭环优化；直接确认；自适应大邻域搜索；空地协同